Website fingerprinting (WF) infers visited websites from observable patterns in encrypted Tor traffic. Although modern classifiers can recognize traffic collected under conditions resembling those of their training data, temporal changes in website content and collection environments can substantially impair generalization. We investigate whether longitudinal traffic information, used alongside a conventional within-visit fingerprint, can improve classification under temporal distribution shift. We evaluate \LTD{}, a model family comprising within-trace (micro), longitudinal (macro), and hybrid configurations, under a chronological training protocol that restricts fitting and model selection to historical development data. We report results on two previously inspected, retrospective evaluation cohorts, covering 65 websites: 09B (14,687 traces) and Future-B (18,543 traces). The hybrid attains 97.440\% and 49.453\% accuracy, respectively, compared with 95.002\% and 43.116\% for its micro-only configuration. On Future-B, it also exceeds Deep Fingerprinting (37.464\%), VarCNN (33.447\%), a CUMUL-style SVM (7.733\%), and a WFlib-based Tik-Tok implementation (20.105\%). Paired block-day analyses of the DF and VarCNN comparisons yield exploratory evidence of larger Future-B accuracy, while additional comparisons and the contrastive NetCLR baseline remain pending. Since evaluation cohorts were examined during method development, the results should not be interpreted as prospective or causally isolating concept drift. The findings motivate longitudinally informed WF as a candidate for improving performance under temporal shifts, while highlighting the need for independent prospective validation.
